% GENERATED FILE. Edit canonical lessons, then run npm run generate:books.
% canonical-source-hash: fnv1a64:45e6bf0f8d5fe487
% canonical-lessons: JA-C83-kono, JA-C83-ano, JA-C83-dono, JA-C83-konna, JA-C83-anna

\chapter{This, That, Which, This kind of, That kind of}
\label{ch:ja-this-that-which-this-kind-of-that-kind-of}
\hlchaptermodality{83}

\begin{chapteropening}
\textbf{By the end of this chapter:} \emph{I can say this, that, which, this kind of and that kind of.}

Complete the last lesson of chapter 83: i can say this, that, which, this kind of and that kind of.
\end{chapteropening}

\section[\texorpdfstring{\ja{この} (kono)}{この (kono)}]{\ja{この} (kono) — this (next to me)}

\label{lesson:JA-C83-kono}



\begin{quote}
Before the new one: say the Japanese for the opposite side, then the Japanese for upstream.
\end{quote}

\subsection*{You'll want to know: \ja{この}}
\textbf{\ja{この}} — \emph{kono} — \textquotedblleft{}this (next to me)\textquotedblright{}.

\begin{grammarlens}[title={What you've built}]
One more word for this chapter.
\end{grammarlens}

\subsection*{Your turn}
\begin{itemize}
\raggedright
  \item \emph{Say it:} \emph{kono}
  \item \emph{Say it:} \emph{kono}, once more
  \item \emph{Say it:} say \emph{kono}
  \item \emph{Recall:} say the Japanese for the opposite side, then the Japanese for upstream, then say \emph{kono} again
\end{itemize}

\begin{tcolorbox}[breakable,colback=teal!4,colframe=teal!35!black,title={Before you move on}]
What does \ja{この} mean? (\textquotedblleft{}This (next to me)\textquotedblright{}.) Say it once more.
\end{tcolorbox}
\section[\texorpdfstring{\ja{あの} (ano)}{あの (ano)}]{\ja{あの} (ano) — that (over there)}

\label{lesson:JA-C83-ano}



\begin{quote}
Before the new one: say the Japanese for upstream, then the Japanese for this.
\end{quote}

\subsection*{You'll want to know: \ja{あの}}
\textbf{\ja{あの}} — \emph{ano} — \textquotedblleft{}that (over there)\textquotedblright{}.

\begin{grammarlens}[title={What you've built}]
One more word for this chapter.
\end{grammarlens}

\subsection*{Your turn}
\begin{itemize}
\raggedright
  \item \emph{Say it:} \emph{ano}
  \item \emph{Say it:} \emph{ano}, once more
  \item \emph{Say it:} say \emph{ano}
  \item \emph{Recall:} say the Japanese for upstream, then the Japanese for this, then say \emph{ano} again
\end{itemize}

\begin{tcolorbox}[breakable,colback=teal!4,colframe=teal!35!black,title={Before you move on}]
What does \ja{あの} mean? (\textquotedblleft{}That (over there)\textquotedblright{}.) Say it once more.
\end{tcolorbox}
\section[\texorpdfstring{\ja{どの} (dono)}{どの (dono)}]{\ja{どの} (dono) — which}

\label{lesson:JA-C83-dono}



\begin{quote}
Before the new one: say the Japanese for this, then the Japanese for that.
\end{quote}

\subsection*{You'll want to know: \ja{どの}}
\textbf{\ja{どの}} — \emph{dono} — \textquotedblleft{}which\textquotedblright{}.

\begin{grammarlens}[title={What you've built}]
One more word for this chapter.
\end{grammarlens}

\subsection*{Your turn}
\begin{itemize}
\raggedright
  \item \emph{Say it:} \emph{dono}
  \item \emph{Say it:} \emph{dono}, once more
  \item \emph{Say it:} say \emph{dono}
  \item \emph{Recall:} say the Japanese for this, then the Japanese for that, then say \emph{dono} again
\end{itemize}

\begin{tcolorbox}[breakable,colback=teal!4,colframe=teal!35!black,title={Before you move on}]
What does \ja{どの} mean? (\textquotedblleft{}Which\textquotedblright{}.) Say it once more.
\end{tcolorbox}
\section[\texorpdfstring{\ja{こんな} (konna)}{こんな (konna)}]{\ja{こんな} (konna) — this kind of}

\label{lesson:JA-C83-konna}



\begin{quote}
Before the new one: say the Japanese for that, then the Japanese for which.
\end{quote}

\subsection*{You'll want to know: \ja{こんな}}
\textbf{\ja{こんな}} — \emph{konna} — \textquotedblleft{}this kind of\textquotedblright{}.

\begin{grammarlens}[title={What you've built}]
One more word for this chapter.
\end{grammarlens}

\subsection*{Your turn}
\begin{itemize}
\raggedright
  \item \emph{Say it:} \emph{konna}
  \item \emph{Say it:} \emph{konna}, once more
  \item \emph{Say it:} say \emph{konna}
  \item \emph{Recall:} say the Japanese for that, then the Japanese for which, then say \emph{konna} again
\end{itemize}

\begin{tcolorbox}[breakable,colback=teal!4,colframe=teal!35!black,title={Before you move on}]
What does \ja{こんな} mean? (\textquotedblleft{}This kind of\textquotedblright{}.) Say it once more.
\end{tcolorbox}
\section[\texorpdfstring{\ja{あんな} (anna)}{あんな (anna)}]{\ja{あんな} (anna) — that kind of}

\label{lesson:JA-C83-anna}



\begin{quote}
Before the new one: say the Japanese for which, then the Japanese for this kind of.

Count long thin things: one, three, six. (\textbf{\ja{いっぽん}}, \textbf{\ja{さんぼん}}, \textbf{\ja{ろっぽん}} — the counter bends after a held beat and after \textbf{\ja{ん}}.)
\end{quote}

\subsection*{You'll want to know: \ja{あんな}}
\textbf{\ja{あんな}} — \emph{anna} — \textquotedblleft{}that kind of\textquotedblright{}.

\begin{grammarlens}[title={What you've built}]
One more word for this chapter.
\end{grammarlens}

\subsection*{Your turn}
\begin{itemize}
\raggedright
  \item \emph{Say it:} \emph{anna}
  \item \emph{Say it:} \emph{anna}, once more
  \item \emph{Say it:} say \emph{anna}
  \item \emph{Recall:} say the Japanese for which, then the Japanese for this kind of, then say \emph{anna} again
\end{itemize}

\begin{tcolorbox}[breakable,colback=teal!4,colframe=teal!35!black,title={Before you move on}]
What does \ja{あんな} mean? (\textquotedblleft{}That kind of\textquotedblright{}.) Say it once more.
\end{tcolorbox}
